\section{Union-Find / Disjoint Set Union}

The \textbf{Union-Find} (or \textbf{Disjoint Set Union}, DSU) data structure maintains a collection
of disjoint sets and supports two primary operations:

\begin{itemize}
\item \textbf{Find}$(x)$: Determine which set element $x$ belongs to (returns the root representative).
\item \textbf{Union}$(x, y)$: Merge the sets containing $x$ and $y$.
\end{itemize}

With \textit{path compression} and \textit{union by rank}, both operations achieve near-constant
amortized time:

$$\alpha(n) \approx O(1) \quad \text{(inverse Ackermann function)}$$

The parent array encodes a forest of trees. For any element $x$:

$$\text{Find}(x) = \begin{cases} x & \text{if } parent[x] = x \\ \text{Find}(parent[x]) & \text{otherwise} \end{cases}$$

Union by rank ensures the shorter tree is attached under the taller one:

$$\text{Union}(x, y): \quad \text{if } rank[r_x] < rank[r_y], \; parent[r_x] \leftarrow r_y$$

\begin{diagram}[id="uf-initial-graph", label="Initial graph with 8 nodes"]
\shape{G}{Graph}{nodes=["0","1","2","3","4","5","6","7"], edges=[("0","1"),("1","2"),("3","4"),("5","6"),("6","7"),("2","3"),("4","5")], directed=false, layout="stable", layout_seed=42}
\recolor{G.node[0]}{state=good}
\recolor{G.node[7]}{state=good}
\annotate{G.node[0]}{label="start", position=above, color=info}
\annotate{G.node[7]}{label="end", position=above, color=info}
\end{diagram}

\subsection{Step-by-Step Animation}

We process edges one at a time, performing Union operations. The animation shows
the graph, the $parent[]$ array, the $rank[]$ array, and a variable watch panel
tracking the current operation.

\begin{animation}[id="uf-walkthrough", label="Union-Find with Path Compression"]
\compute{
initial_parent = [0, 1, 2, 3, 4, 5, 6, 7]
initial_rank = [0, 0, 0, 0, 0, 0, 0, 0]
all_nodes = [0, 1, 2, 3, 4, 5, 6, 7]
edges_to_process = ["(0,1)", "(1,2)", "(3,4)", "(5,6)", "(6,7)", "(2,3)", "(4,5)"]
}

\shape{G}{Graph}{nodes=["0","1","2","3","4","5","6","7"], edges=[("0","1"),("1","2"),("3","4"),("5","6"),("6","7"),("2","3"),("4","5")], directed=false, layout="stable", layout_seed=42}
\shape{par}{Array}{size=8, data=[0,1,2,3,4,5,6,7], labels="0..7", label="$parent[]$"}
\shape{rnk}{Array}{size=8, data=[0,0,0,0,0,0,0,0], labels="0..7", label="$rank[]$"}
\shape{vars}{VariableWatch}{names=["op","x","y","root_x","root_y","step"], label="Variables"}

% Step 1: Initial state
\step
\apply{vars.var[op]}{value="init"}
\apply{vars.var[x]}{value="-"}
\apply{vars.var[y]}{value="-"}
\apply{vars.var[root_x]}{value="-"}
\apply{vars.var[root_y]}{value="-"}
\apply{vars.var[step]}{value="0"}
\narrate{Initial state: 8 disjoint elements $\{0\}, \{1\}, \ldots, \{7\}$. Each element is its own parent, all ranks are 0.}

% Step 2: Union(0, 1)
\step
\apply{vars.var[op]}{value="Union(0,1)"}
\apply{vars.var[x]}{value="0"}
\apply{vars.var[y]}{value="1"}
\apply{vars.var[root_x]}{value="0"}
\apply{vars.var[root_y]}{value="1"}
\apply{vars.var[step]}{value="1"}
\recolor{G.node[0]}{state=current}
\recolor{G.node[1]}{state=current}
\recolor{G.edge[(0,1)]}{state=current}
\highlight{par.cell[1]}
\narrate{Union(0, 1): Both roots have rank 0. Attach 1 under 0 and increment rank of 0.}

% Step 3: Apply Union(0, 1)
\step
\apply{par.cell[1]}{value=0}
\apply{rnk.cell[0]}{value=1}
\recolor{par.cell[1]}{state=done}
\recolor{rnk.cell[0]}{state=done}
\recolor{G.node[0]}{state=good}
\recolor{G.node[1]}{state=done}
\recolor{G.edge[(0,1)]}{state=good}
\annotate{par.cell[1]}{label="0 is root", position=below, color=good, arrow_from="par.cell[0]"}
\narrate{After Union(0,1): $parent[1] = 0$, $rank[0] = 1$. Set $\{0, 1\}$ formed with root 0.}

% Step 4: Union(1, 2)
\step
\apply{vars.var[op]}{value="Union(1,2)"}
\apply{vars.var[x]}{value="1"}
\apply{vars.var[y]}{value="2"}
\apply{vars.var[root_x]}{value="0"}
\apply{vars.var[root_y]}{value="2"}
\apply{vars.var[step]}{value="2"}
\recolor{G.node[2]}{state=current}
\recolor{G.edge[(1,2)]}{state=current}
\highlight{par.cell[2]}
\narrate{Union(1, 2): Find(1) = 0 (rank 1), Find(2) = 2 (rank 0). Attach 2 under 0.}

% Step 5: Apply Union(1, 2)
\step
\apply{par.cell[2]}{value=0}
\recolor{par.cell[2]}{state=done}
\recolor{G.node[2]}{state=done}
\recolor{G.edge[(1,2)]}{state=good}
\narrate{After Union(1,2): $parent[2] = 0$. Set $\{0, 1, 2\}$ with root 0.}

% Step 6: Union(3, 4)
\step
\apply{vars.var[op]}{value="Union(3,4)"}
\apply{vars.var[x]}{value="3"}
\apply{vars.var[y]}{value="4"}
\apply{vars.var[root_x]}{value="3"}
\apply{vars.var[root_y]}{value="4"}
\apply{vars.var[step]}{value="3"}
\recolor{G.node[3]}{state=current}
\recolor{G.node[4]}{state=current}
\recolor{G.edge[(3,4)]}{state=current}
\highlight{par.cell[4]}
\narrate{Union(3, 4): Both roots have rank 0. Attach 4 under 3 and increment rank of 3.}

% Step 7: Apply Union(3, 4)
\step
\apply{par.cell[4]}{value=3}
\apply{rnk.cell[3]}{value=1}
\recolor{par.cell[4]}{state=done}
\recolor{rnk.cell[3]}{state=done}
\recolor{G.node[3]}{state=good}
\recolor{G.node[4]}{state=done}
\recolor{G.edge[(3,4)]}{state=good}
\annotate{par.cell[4]}{label="3 is root", position=below, color=good, arrow_from="par.cell[3]"}
\narrate{After Union(3,4): $parent[4] = 3$, $rank[3] = 1$. Set $\{3, 4\}$ formed with root 3.}

% Step 8: Union(5, 6)
\step
\apply{vars.var[op]}{value="Union(5,6)"}
\apply{vars.var[x]}{value="5"}
\apply{vars.var[y]}{value="6"}
\apply{vars.var[root_x]}{value="5"}
\apply{vars.var[root_y]}{value="6"}
\apply{vars.var[step]}{value="4"}
\recolor{G.node[5]}{state=current}
\recolor{G.node[6]}{state=current}
\recolor{G.edge[(5,6)]}{state=current}
\highlight{par.cell[6]}
\narrate{Union(5, 6): Both roots have rank 0. Attach 6 under 5 and increment rank of 5.}

% Step 9: Apply Union(5, 6)
\step
\apply{par.cell[6]}{value=5}
\apply{rnk.cell[5]}{value=1}
\recolor{par.cell[6]}{state=done}
\recolor{rnk.cell[5]}{state=done}
\recolor{G.node[5]}{state=good}
\recolor{G.node[6]}{state=done}
\recolor{G.edge[(5,6)]}{state=good}
\narrate{After Union(5,6): $parent[6] = 5$, $rank[5] = 1$. Set $\{5, 6\}$ formed with root 5.}

% Step 10: Union(6, 7) -- Find(6) triggers path to 5
\step
\apply{vars.var[op]}{value="Union(6,7)"}
\apply{vars.var[x]}{value="6"}
\apply{vars.var[y]}{value="7"}
\apply{vars.var[root_x]}{value="5"}
\apply{vars.var[root_y]}{value="7"}
\apply{vars.var[step]}{value="5"}
\recolor{G.node[7]}{state=current}
\recolor{G.edge[(6,7)]}{state=current}
\highlight{par.cell[7]}
\narrate{Union(6, 7): Find(6) = 5 (rank 1), Find(7) = 7 (rank 0). Attach 7 under 5.}

% Step 11: Apply Union(6, 7)
\step
\apply{par.cell[7]}{value=5}
\recolor{par.cell[7]}{state=done}
\recolor{G.node[7]}{state=done}
\recolor{G.edge[(6,7)]}{state=good}
\narrate{After Union(6,7): $parent[7] = 5$. Set $\{5, 6, 7\}$ with root 5.}

% Step 12: Union(2, 3) -- merges two larger sets
\step
\apply{vars.var[op]}{value="Union(2,3)"}
\apply{vars.var[x]}{value="2"}
\apply{vars.var[y]}{value="3"}
\apply{vars.var[root_x]}{value="0"}
\apply{vars.var[root_y]}{value="3"}
\apply{vars.var[step]}{value="6"}
\recolor{G.edge[(2,3)]}{state=current}
\recolor{G.node[0]}{state=current}
\recolor{G.node[3]}{state=current}
\highlight{par.cell[3]}
\annotate{par.cell[3]}{label="rank tie!", position=above, color=warn}
\narrate{Union(2, 3): Find(2) = 0 (rank 1), Find(3) = 3 (rank 1). Equal ranks! Attach 3 under 0, increment rank of 0.}

% Step 13: Apply Union(2, 3)
\step
\apply{par.cell[3]}{value=0}
\apply{rnk.cell[0]}{value=2}
\recolor{par.cell[3]}{state=done}
\recolor{rnk.cell[0]}{state=done}
\recolor{G.edge[(2,3)]}{state=good}
\recolor{G.node[0]}{state=good}
\recolor{G.node[3]}{state=done}
\narrate{After Union(2,3): $parent[3] = 0$, $rank[0] = 2$. Set $\{0,1,2,3,4\}$ with root 0.}

% Step 14: Union(4, 5) -- merges everything
\step
\apply{vars.var[op]}{value="Union(4,5)"}
\apply{vars.var[x]}{value="4"}
\apply{vars.var[y]}{value="5"}
\apply{vars.var[root_x]}{value="0"}
\apply{vars.var[root_y]}{value="5"}
\apply{vars.var[step]}{value="7"}
\recolor{G.edge[(4,5)]}{state=current}
\highlight{par.cell[5]}
\narrate{Union(4, 5): Find(4) $\rightarrow$ 3 $\rightarrow$ 0 (rank 2), Find(5) = 5 (rank 1). Attach 5 under 0. Path compression on Find(4)!}

% Step 15: Apply Union(4, 5) with path compression
\step
\apply{par.cell[5]}{value=0}
\apply{par.cell[4]}{value=0}
\recolor{par.cell[5]}{state=done}
\recolor{par.cell[4]}{state=path}
\recolor{G.edge[(4,5)]}{state=good}
\recolor{G.node[5]}{state=done}
\annotate{par.cell[4]}{label="compressed!", position=above, color=path}
\narrate{After Union(4,5): $parent[5] = 0$. Path compression: $parent[4]$ updated from 3 to 0 directly. All 8 elements now in one set!}

% Step 16: Find(7) with path compression
\step
\apply{vars.var[op]}{value="Find(7)"}
\apply{vars.var[x]}{value="7"}
\apply{vars.var[y]}{value="-"}
\apply{vars.var[root_x]}{value="0"}
\apply{vars.var[root_y]}{value="-"}
\apply{vars.var[step]}{value="8"}
\recolor{G.node[7]}{state=current}
\recolor{G.node[5]}{state=current}
\recolor{G.node[0]}{state=current}
\narrate{Find(7): traverse $7 \rightarrow 5 \rightarrow 0$. Root is 0. Now apply path compression.}

% Step 17: Path compression result
\step
\apply{par.cell[7]}{value=0}
\recolor{par.cell[7]}{state=path}
\recolor{G.node[7]}{state=done}
\recolor{G.node[5]}{state=done}
\recolor{G.node[0]}{state=good}
\annotate{par.cell[7]}{label="compressed!", position=above, color=path}
\narrate{Path compression: $parent[7]$ updated from 5 to 0. Future Find(7) returns in $O(1)$.}

% Step 18: Final state
\step
\apply{vars.var[op]}{value="done"}
\apply{vars.var[step]}{value="done"}
\compute{
done_nodes = [0, 1, 2, 3, 4, 5, 6, 7]
}
\foreach{i}{${done_nodes}}
  \recolor{par.cell[${i}]}{state=good}
\endforeach
\foreach{i}{${done_nodes}}
  \recolor{rnk.cell[${i}]}{state=good}
\endforeach
\narrate{Final state: all 8 nodes belong to one set with root 0. The $parent[]$ array is nearly flat thanks to path compression.}

\end{animation}

\subsection{Complexity Analysis}

\begin{tabular}{|l|c|c|}
\hline
\textbf{Operation} & \textbf{Naive} & \textbf{With Optimizations} \\
\hline
Find & $O(n)$ & $O(\alpha(n))$ \\
\hline
Union & $O(n)$ & $O(\alpha(n))$ \\
\hline
$m$ operations on $n$ elements & $O(mn)$ & $O(m \cdot \alpha(n))$ \\
\hline
Space & $O(n)$ & $O(n)$ \\
\hline
\end{tabular}

Where $\alpha(n)$ is the \textbf{inverse Ackermann function}, which grows so slowly that
$\alpha(n) \leq 4$ for all practical values of $n$ (up to $2^{2^{2^{65536}}}$).

\textbf{Key optimizations:}
\begin{enumerate}
\item \textbf{Path compression}: During Find, make every node on the path point directly to the root.
\item \textbf{Union by rank}: Always attach the shorter tree under the taller one to keep trees balanced.
\end{enumerate}

Together, these guarantee that any sequence of $m$ Union/Find operations on $n$ elements
takes $O(m \cdot \alpha(n))$ time, which is effectively $O(m)$ in practice.
